%! Matrix Computations and A = LU — Unit 2, Lesson 2.3 — Slides (Beamer)
\documentclass[aspectratio=169, 11pt]{beamer}
\usepackage{linalg-beamer}   % provides \burgundyheader{} and \sectionlabel[color]{}
\newcommand{\xx}{\mathbf{x}}
\newcommand{\bb}{\mathbf{b}}
\newcommand{\cc}{\mathbf{c}}

\begin{document}

% TITLE SLIDE (CourseName is not defined in beamer, so the name is literal)
{
\setbeamercolor{background canvas}{bg=burgundy}
\begin{frame}[plain]
  \vspace{2.8cm}\hspace{0.55cm}%
  \begin{minipage}{0.9\textwidth}
    {\color{goldacc}\bfseries\small LESSON 2.3}\\[0.5cm]
    {\color{white}\bfseries\LARGE Matrix Computations and $A=LU$}\\[0.3cm]
    {\color{white}\bfseries\large Elimination, saved --- factor once, solve fast}\\[1.0cm]
    {\color{blushmid}\small Linear Algebra~$\cdot$~Unit 2: Solving Linear Equations $A\xx=\bb$}
  \end{minipage}
\end{frame}
}

% HOOK
\begin{frame}[t, plain]
  \burgundyheader{Same recipe, a new order every week}
  \sectionlabel{Hook}
  Last lesson's bakery: the recipe matrix never changes --- only the order $\bb$ does.
  \[
    A = \begin{bmatrix} 1 & 3 \\ 2 & 7 \end{bmatrix}, \qquad
    \text{this week } \bb = \begin{bmatrix} 7 \\ 15 \end{bmatrix}.
  \]
  \begin{block}{The question}
    In 2.2 we built the whole inverse $A^{-1}$. But elimination in 2.1 already did the hard work ---
    can we just \emph{save} it and reuse it for every order?
  \end{block}
\end{frame}

% SAVE ELIMINATION AS LU
\begin{frame}[t, plain]
  \burgundyheader{Save elimination as $A=LU$}
  \sectionlabel[goldacc]{The idea}
  \begin{itemize}
    \setlength{\itemsep}{5pt}
    \item Elimination ($\ell=2$) turns $A$ into the upper-triangular $U=\begin{bsmallmatrix} 1&3\\0&1 \end{bsmallmatrix}$.
    \item Put the multiplier back (with a $+$) to get the lower-triangular $L=\begin{bsmallmatrix} 1&0\\2&1 \end{bsmallmatrix}$.
  \end{itemize}
  \begin{block}{$LU$ rebuilds $A$}
    $\begin{bsmallmatrix} 1&0\\2&1 \end{bsmallmatrix}\begin{bsmallmatrix} 1&3\\0&1 \end{bsmallmatrix} = \begin{bsmallmatrix} 1&3\\2&7 \end{bsmallmatrix} = A.$ \quad
    $U$ = \emph{what} elimination got; $L$ = \emph{how} we did it.
  \end{block}
\end{frame}

% L STORES THE MULTIPLIERS
\begin{frame}[t, plain]
  \burgundyheader{$L$ just stores the multipliers}
  \sectionlabel{No extra work}
  Eliminate $A=\begin{bsmallmatrix} 1&2&1\\2&5&5\\3&7&8 \end{bsmallmatrix}$ and record each multiplier:
  \; $\ell_{21}=2,\ \ell_{31}=3,\ \ell_{32}=1$.
  \[
    U=\begin{bmatrix} 1&2&1\\0&1&3\\0&0&2 \end{bmatrix},
    \qquad
    L=\begin{bmatrix} 1&0&0\\ 2&1&0\\ 3&1&1 \end{bmatrix}.
  \]
  The multipliers slot straight into $L$ --- it is a \textbf{free record} of the elimination you already did.
\end{frame}

% TWO SWEEPS
\begin{frame}[t, plain]
  \burgundyheader{Solve in two triangular sweeps}
  \sectionlabel[goldacc]{Factor once, solve fast}
  $A\xx=\bb$ becomes $L(U\xx)=\bb$. Name $\cc=U\xx$, then solve two easy systems:
  \[
    \underbrace{\begin{bmatrix} 1&0\\2&1 \end{bmatrix}\cc = \begin{bmatrix} 7\\15 \end{bmatrix}}_{\text{forward: }L\cc=\bb}
    \Rightarrow \cc=\begin{bmatrix} 7\\1 \end{bmatrix},
    \quad
    \underbrace{\begin{bmatrix} 1&3\\0&1 \end{bmatrix}\xx = \begin{bmatrix} 7\\1 \end{bmatrix}}_{\text{back: }U\xx=\cc}
    \Rightarrow \xx=\begin{bmatrix} 4\\1 \end{bmatrix}.
  \]
  The forward sweep \emph{is} ``update $\bb$'' (note $\cc=\begin{bsmallmatrix} 7\\1 \end{bsmallmatrix}$ from 2.1); the back sweep is old back-substitution.
\end{frame}

% WHY BOTHER
\begin{frame}[t, plain]
  \burgundyheader{Why factor? The computation cost}
  \sectionlabel{Matrix computations}
  For one $A$ and \emph{many} right-hand sides $\bb$:
  \begin{itemize}
    \setlength{\itemsep}{5pt}
    \item Re-run elimination each time --- repeats all the hard work every order.
    \item Build $A^{-1}$ --- powerful, but more work to build than $L$ and $U$.
    \item \textbf{Factor once, two sweeps per order} --- cheapest to solve.
  \end{itemize}
  \begin{block}{The takeaway}
    Factoring is the expensive step, done \emph{once}; each new order is two quick sweeps. This is how software actually solves $A\xx=\bb$.
  \end{block}
\end{frame}

% PREVIEW
\begin{frame}[t, plain]
  \burgundyheader{Where this is heading}
  \sectionlabel{Connections}
  \textbf{Today:} elimination $\to$ save as $A=LU$ $\to$ solve in two sweeps ($L\cc=\bb$, $U\xx=\cc$).\\[8pt]
  \begin{itemize}
    \setlength{\itemsep}{4pt}
    \item \textbf{2.4} Permutations and Transposes --- if a pivot is $0$, swap rows: $PA=LU$.
    \item Beyond: $A=LU$ is the workhorse behind fast, repeated solving.
  \end{itemize}
\end{frame}

\end{document}
